% !TEX root = 00_dmlr_compel.tex
\section{Related Work}\label{sec:related}

Data selection pipelines for pretraining large models typically involve a sequence of filtering stages to manage the scale and heterogeneity of web-derived corpora.
Common goals include removing undesirable content---such as duplicated boilerplate, SEO spam, or template-heavy pages---and maximizing overall data quality.

\paragraph{Language filtering.}
Language identification is usually performed using lightweight classifiers such as fastText or langdetect \citep{raffel2020exploring, soldaini2024dolma, conneau2019cross, xue2021mt5, laurencon2022roots, grave2018learning}.
These methods are efficient but require careful threshold tuning to avoid excessive false positives or negatives.

\paragraph{Heuristic filtering.}
Surface-level heuristics such as document length or the presence of blacklisted terms are commonly used \citep{rae2021scaling, raffel2020exploring, aghajanyan2022htlm}.
These are computationally cheap but often fail to capture more nuanced indicators of quality.

\paragraph{Quality filtering.}
Retaining educational, human-written text is a key objective in many pipelines.
Classifier-based approaches typically train fastText models or hashed-feature classifiers on curated datasets \citep{du2022glam, longpre2023flan}, while perplexity-based methods score documents using language models trained on high-quality corpora \citep{wenzek2020ccnet, wettig2024qurating}.
Despite strong performance, these methods face limitations: (1) reference corpora may encode cultural or demographic preferences, introducing bias that can marginalize certain dialects, domains, or communities \citep{gururangan2022whose, xu2021detoxifying}, thereby reinforcing existing inequities in model behavior, and (2) the notion of ``quality'' remains inherently subjective and task-dependent \citep{longpre2023flan}.
In contrast, \textsc{Compel} sidesteps these challenges by avoiding reliance on human-annotated or culturally-specific reference corpora, instead using a purely statistical signal (compression ratio) that generalizes across domains.

\paragraph{LM-based quality filtering.}
Recent approaches have employed language models for data refinement.
For instance, \citet{maini2024rephrasing} introduced Web Rephrase Augmented Pre-training (WRAP), which uses instruction-tuned models to paraphrase web documents into styles like Wikipedia or question-answer formats, enhancing data quality and training efficiency.
Similarly, \citet{su2024nemotron} presented Nemotron-CC, a refined dataset derived from Common Crawl, combining classifier ensembling, synthetic data rephrasing, and reduced reliance on heuristic filters to balance data quality and quantity.
While effective, these methods are computationally intensive, requiring substantial resources for model inference and data generation.

\paragraph{Compression-based data selection.}
The link between data compression and language modeling has long been observed \citep{shannon1951prediction}.
Recent studies reinforce this connection: LMs can act as powerful compressors \citep{deletang2023language}, and generalization ability may correlate with compression capacity \citep{huang2024compression}.
\citet{pandey2024gzip} proposes a data-dependent scaling law based on gzip compressibility.
Compression has recently emerged as a scalable, embedding-free signal for data selection and quality estimation.
\citet{yin2024entropy} selects diverse alignment data by scoring examples based on their incremental compressed size, optimizing corpus diversity.
ZIP-FIT \citep{obbad2024zipfit} uses Normalized Compression Distance (NCD) to filter source data for fine-tuning tasks, outperforming neural embeddings in settings such as code and autoformalization.

\textsc{Compel} differs from ZIP and ZIP-FIT in both goal and design.
ZIP aims to construct diverse, low-redundancy subsets for alignment by selecting examples with low incremental compressed size, relying on early-stage training signals and task-specific assumptions.
ZIP-FIT, in contrast, optimizes similarity to a target task distribution using pairwise Normalized Compression Distance.
\textsc{Compel}, by comparison, is a general-purpose filter designed for broad pretraining use: it discards overly redundant or noisy documents based purely on raw compression ratio.
It requires no task-specific supervision, no pairwise comparisons, and no access to model loss signals.
Its simplicity and speed make it especially attractive for trillion-token pipelines.
